% ---------------------------------------------------------------
% TEMPLATE TCC1 - SENAC SANTO AMARO
% ---------------------------------------------------------------
% Trabalho de Conclusão de Curso 1
% Elaborado para apresentação da proposta inicial do TCC
% Estrutura: 4 capítulos (Introdução, Revisão Bibliográfica,
%            Desenvolvimento, Referências Bibliográficas)
% ---------------------------------------------------------------

\documentclass[
    12pt,                % tamanho da fonte
    openright,           % capítulos começam em página ímpar
    oneside,             % para impressão em apenas um lado do papel
    a4paper,             % tamanho do papel
    english,             % idioma adicional
    brazil               % idioma principal
]{abntex2}

% ---------------------------------------------------------------
% PACOTES
% ---------------------------------------------------------------
\usepackage{lmodern}                % Fonte Latin Modern
\usepackage[T1]{fontenc}            % Seleção de códigos de fonte
\usepackage[utf8]{inputenc}         % Codificação do documento
\usepackage{lastpage}               % Usado pela Ficha catalográfica
\usepackage{indentfirst}            % Indenta o primeiro parágrafo
\usepackage{color}                  % Controle de cores
\usepackage{graphicx}               % Inclusão de gráficos
\usepackage{microtype}              % Melhorias de justificação
\usepackage{float}                  % Para posicionamento de figuras
\usepackage{booktabs}               % Tabelas profissionais
\usepackage{multirow}               % Células mescladas em tabelas
\usepackage{longtable}              % Tabelas longas
\usepackage{amsmath,amssymb}        % Símbolos matemáticos
\usepackage{pgfgantt}               % Diagramas de Gantt
\usepackage[brazilian,hyperpageref]{backref}  % Páginas com citações
\usepackage[alf]{abntex2cite}       % Citações ABNT

% ---------------------------------------------------------------
% CONFIGURAÇÕES DO PDF
% ---------------------------------------------------------------
\makeatletter
\hypersetup{
    pdftitle={\@title},
    pdfauthor={\@author},
    pdfsubject={TCC1 - SENAC Santo Amaro},
    pdfcreator={LaTeX with abnTeX2},
    pdfkeywords={tcc}{senac}{santo amaro}{trabalho de conclusão},
    colorlinks=true,        % false: links em caixas; true: links coloridos
    linkcolor=blue,         % cor dos links internos
    citecolor=blue,         % cor dos links para bibliografia
    filecolor=magenta,      % cor dos links para arquivos
    urlcolor=blue,
    bookmarksdepth=4
}
\makeatother

% ---------------------------------------------------------------
% CONFIGURAÇÕES DE APARÊNCIA DO PDF FINAL
% ---------------------------------------------------------------
\definecolor{blue}{RGB}{41,5,195}

% ---------------------------------------------------------------
% INFORMAÇÕES DO DOCUMENTO
% ---------------------------------------------------------------
\titulo{Detecção de Intenção Implícita para Ativação de Assistentes de Voz em Ambientes Controlados}
\autor{Phelipe Pereira de Souza}
\local{São Paulo}
\data{2026}
\orientador{Prof. Thyago Conchado Quintas}
% \coorientador{Prof. Nome do Coorientador} % Descomente se houver coorientador

\instituicao{%
  SENAC -- Serviço Nacional de Aprendizagem Comercial
  \par
  Unidade Santo Amaro
  \par
  Bacharelado em Ciência da Computação
}

\tipotrabalho{Trabalho de Conclusão de Curso I}

\preambulo{Proposta de Trabalho de Conclusão de Curso apresentado ao SENAC Santo Amaro como requisito parcial para aprovação na disciplina de TCC1 do Bacharelado em Ciência da Computação.}

% ---------------------------------------------------------------
% COMPILA O ÍNDICE
% ---------------------------------------------------------------
\makeindex

% ---------------------------------------------------------------
% INÍCIO DO DOCUMENTO
% ---------------------------------------------------------------
\begin{document}

% Seleciona o idioma do documento
\selectlanguage{brazil}

% Retira espaço extra obsoleto entre as frases
\frenchspacing

% ---------------------------------------------------------------
% ELEMENTOS PRÉ-TEXTUAIS
% ---------------------------------------------------------------
\pretextual

% ---------------------------------------------------------------
% CAPA
% ---------------------------------------------------------------
\imprimircapa

% ---------------------------------------------------------------
% FOLHA DE ROSTO
% ---------------------------------------------------------------
\imprimirfolhaderosto

% ---------------------------------------------------------------
% RESUMO
% ---------------------------------------------------------------
% \begin{resumo}

    
% \vspace{\onelineskip}

% \noindent
% % \textbf{Palavras-chave}: palavra1. palavra2. palavra3. palavra4. palavra5.
% \end{resumo}

% ---------------------------------------------------------------
% LISTA DE ILUSTRAÇÕES
% ---------------------------------------------------------------
% \pdfbookmark[0]{\listfigurename}{lof}
% \listoffigures*
% \cleardoublepage

% ---------------------------------------------------------------
% LISTA DE TABELAS
% ---------------------------------------------------------------
% \pdfbookmark[0]{\listtablename}{lot}
% \listoftables*
% \cleardoublepage

% ---------------------------------------------------------------
% SUMÁRIO
% ---------------------------------------------------------------
\pdfbookmark[0]{\contentsname}{toc}
\tableofcontents*
\cleardoublepage

% ---------------------------------------------------------------
% ELEMENTOS TEXTUAIS
% ---------------------------------------------------------------
\textual

% ===============================================================
% CAPÍTULO 1 - INTRODUÇÃO
% ===============================================================
\chapter{Introdução}
\label{cap:introducao}

Nas últimas décadas, os assistentes virtuais baseados em voz tornaram-se componentes cada vez mais presentes no cotidiano, sendo incorporados em \textit{smartphones}, alto-falantes inteligentes, \textit{wearables} e sistemas de automação residencial \cite{Jurafsky2024}. Esses sistemas permitem que usuários interajam com dispositivos computacionais por meio de comandos de voz, facilitando a execução de tarefas como busca de informações, controle de dispositivos domésticos e gerenciamento de atividades. O mercado global de assistentes virtuais baseados em voz foi avaliado em US\$ 7,35 bilhões em 2024 e é projetado para alcançar US\$ 33,74 bilhões até 2030, com uma taxa de crescimento anual composta de 26,5\%, evidenciando a escala e a relevância dessa tecnologia \cite{NMSC2025}.

Na maioria dos sistemas atuais, a interação com assistentes virtuais baseados em voz é iniciada por meio de palavras de ativação explícitas, conhecidas como \textit{wake-words} ou \textit{hotwords}, tais como ``Hey Siri'', ``Alexa'' ou ``Ok Google''. Essas expressões atuam como um mecanismo de gatilho que sinaliza ao sistema o desejo do usuário de iniciar uma interação. Do ponto de vista técnico, essa abordagem é eficaz, pois permite reduzir o processamento contínuo de áudio e minimizar ativações indevidas \cite{Garg2021}. No entanto, a dependência de \textit{wake-words} impõe limitações significativas à naturalidade da comunicação humano-computador, tornando a interação menos fluida e intuitiva quando comparada ao diálogo humano \cite{Nie2021}.

Em interações humanas cotidianas, comandos e solicitações são frequentemente expressos de forma implícita, inseridos naturalmente no fluxo da conversação, sem a necessidade de uma palavra específica para iniciar a comunicação \cite{Weld2022}. Em contraste, os sistemas baseados em \textit{wake-words} são vulneráveis ao fenômeno de ativações falsas por palavras foneticamente semelhantes, conhecido como \textit{FakeWake} \cite{Chen2021}. Esse problema compromete a qualidade da experiência, reforçando as limitações da abordagem atual para uma interação verdadeiramente natural e intuitiva.

Nesse cenário, surge a área de pesquisa denominada Detecção de Fala Direcionada ao Dispositivo, em inglês \textit{Device-Directed Speech Detection} (DDSD), que visa determinar automaticamente se uma fala é direcionada a um assistente virtual, mesmo na ausência de uma palavra de ativação explícita, ou se faz parte de uma conversa entre humanos \cite{Mallidi2018, Wagner2024}. Essa abordagem possibilita interações mais naturais e fluidas, reduzindo a necessidade de gatilhos verbais explícitos (\textit{wake-words}) e representando um avanço significativo rumo a uma comunicação humano-computador mais intuitiva e semelhante ao diálogo cotidiano.

Este trabalho investiga a viabilidade da DDSD em português brasileiro por meio de uma abordagem baseada exclusivamente no nível textual. O sistema proposto opera com transcrições de fala obtidas via Reconhecimento Automático de Fala, em inglês \textit{Automatic Speech Recognition} (ASR), e aplica técnicas de Processamento de Linguagem Natural (PLN) e Aprendizado de Máquina, em inglês \textit{Machine Learning} (ML), para classificar automaticamente se uma fala é direcionada ao dispositivo ou faz parte de uma conversa incidental entre humanos. O foco central está na etapa de classificação de intenção; o ASR é tratado como módulo de pré-processamento responsável pela geração das transcrições textuais utilizadas como entrada do classificador.

% ---------------------------------------------------------------

\section{Cenário motivador}
\label{sec:cenario}

Para ilustrar a dificuldade e a relevância do problema, considere o seguinte cenário: em uma sala de estar, duas pessoas conversam enquanto um assistente virtual permanece ligado. Umas delas diz, em tom casual: ``nossa, está escuro aqui''. O sistema não deve necessariamente acionar a iluminação, pois a frase é provavelmente um comentário, não um comando. Momentos depois, a mesma pessoa olha para o dispositivo e diz: ``acende a luz da sala''. Nesse caso, mesmo sem \textit{wake-word}, o sistema deveria identificar a fala como claramente direcionada ao dispositivo.

Esse cenário ilustra três classes de enunciados que o sistema precisa distinguir: (i) comandos diretos implícitos dirigidos ao dispositivo, como ``acende a luz''; (ii) comentários incidentais que não são comandos, como ``nossa, está escuro aqui''; e (iii) casos ambíguos dependentes de contexto, como ``apaga a luz, por favor'', que pode ser dirigido a uma pessoa ou ao assistente. O problema de DDSD consiste precisamente em classificar automaticamente esses enunciados, sem depender de \textit{wake-words} explícitas \cite{Mallidi2018, Wagner2024}.

Esse tipo de sistema tem aplicações relevantes em diferentes contextos: casas inteligentes, veículos com assistentes de voz, dispositivos para pessoas idosas, \textit{smartphones} e \textit{wearables}. Em todos esses cenários, a necessidade de invocar repetidamente o dispositivo com uma \textit{wake-word} cria uma barreira artificial entre o usuário e o sistema, reduzindo a naturalidade e a fluidez da interação \cite{Nie2021}.

% ---------------------------------------------------------------
\section{Contexto}
\label{sec:contexto}

O desenvolvimento dos assistentes virtuais baseados em voz acompanhou os avanços significativos nas áreas de ASR e PLN. Modelos como o Whisper \cite{Radford2023} representam um marco no campo do ASR, possibilitando transcrições de áudio com alta precisão e robustez em múltiplos idiomas. No entanto, sua arquitetura padrão opera em lotes (\textit{batch}), o que dificulta seu uso direto em interações fluidas que exigem processamento contínuo em tempo real (\textit{streaming}). Para contornar essa limitação, trabalhos como o Whispy \cite{Bevilacqua2024} propuseram adaptações arquiteturais para viabilizar o uso do Whisper nesses cenários dinâmicos. Paralelamente, surgiram modelos projetados para transcrição ao vivo com menor custo computacional, como o Moonshine \cite{Jeffries2024}, apresentado como uma alternativa otimizada para tarefas de comando de voz e assistentes embarcados.

No caso específico do português brasileiro, modelos baseados na arquitetura Wav2Vec 2.0 \cite{Grosman2021}, treinados com dados locais, têm apresentado resultados expressivos, especialmente em cenários com menor disponibilidade de dados rotulados, oferecendo alta precisão e melhor adaptação às particularidades linguísticas do idioma.

Simultaneamente, no campo do PLN, a arquitetura Transformer, proposta por \citeonline{Vaswani2017}, revolucionou o processamento de linguagem natural, servindo de base para os modernos \textit{Large Language Models} (LLMs). No cenário nacional, modelos como o BERTimbau \cite{Souza2020} consolidaram o uso de arquiteturas Transformer adaptadas ao português brasileiro em tarefas de compreensão textual, enquanto abordagens mais recentes, como Albertina \cite{Rodrigues2023} e DeBERTinha \cite{Abonizio2023}, buscam maior eficiência e desempenho em tarefas de representação linguística. Paralelamente, LLMs generativos como o Sabiá \cite{Pires2023} expandiram essas capacidades para aplicações conversacionais e de geração de texto, estabelecendo avanços relevantes no processamento de linguagem natural em português.

É importante delimitar o papel de cada família de modelos neste trabalho: os modelos ASR (Whisper, Moonshine, Wav2Vec 2.0) são utilizados exclusivamente como módulos de transcrição de fala para texto, gerando a entrada textual do classificador. Os modelos de linguagem pré-treinados (BERTimbau, Albertina, DeBERTinha) são os candidatos ao papel de classificador de intenção, componente central da proposta. Modelos como o Sabiá são mencionados apenas como parte do referencial teórico sobre o estado da arte em LLMs para o português. A tabela \ref{tab:tecnologias} sintetiza esse mapeamento.

\begin{table}[H]
\caption{Mapeamento de tecnologias citadas no trabalho}
\label{tab:tecnologias}
\begin{tabular}{lll}
\toprule
\textbf{Tecnologia} & \textbf{Papel no trabalho} & \textbf{Será implementada?} \\
\midrule
Whisper             & ASR - geração de transcrição     & Sim (principal) \\
Moonshine           & ASR - alternativa leve           & Opcional \\
Wav2Vec 2.0 PT-BR   & ASR - referência em português    & Opcional \\
BERTimbau           & Classificador textual (principal) & Sim \\
DeBERTinha          & Classificador textual (alternativo) & Opcional \\
Albertina           & Classificador textual (alternativo) & Opcional \\
TF-IDF + SVM/LR     & Baseline clássico                & Sim \\
Sabiá               & Referencial teórico — LLM PT-BR  & Não \\
\bottomrule
\end{tabular}
\fonte{Elaborado pelo autor.}
\end{table}

Pesquisas recentes na área de DDSD têm explorado abordagens semelhantes à proposta neste trabalho, utilizando transcrições de áudio aliadas a técnicas de PLN para superar as limitações de abordagens exclusivamente acústicas. Trabalhos como os de \citeonline{Wagner2024} e \citeonline{Dighe2024} demonstraram que a aplicação de modelos de linguagem sobre as hipóteses textuais geradas por sistemas ASR é uma estratégia promissora para detectar a intenção de fala direcionada ao dispositivo. Expandindo esse conceito, \citeonline{Rudovic2024} evidenciaram como essa análise baseada em texto permite compreender intenções em interações contínuas e fluxos de diálogo, corroborando a premissa de que a análise textual do áudio transcrito é um caminho viável e promissor.

% ---------------------------------------------------------------
\section{Justificativa}
\label{sec:justificativa}

A crescente presença de assistentes virtuais baseados em voz em diferentes dispositivos computacionais evidencia a relevância de pesquisas voltadas à melhoria de interações humano-computador. Apesar dos avanços significativos em ASR e PLN, a maioria dos sistemas atuais ainda depende de palavras de ativação explícitas (\textit{wake-words}), que podem comprometer a naturalidade e a fluidez da comunicação \cite{Mallidi2018}. Em ambientes dinâmicos, a necessidade de invocar repetidamente o dispositivo pode gerar fadiga no usuário e reduzir a percepção de naturalidade da interação \cite{Nie2021}. Nesse contexto, a DDSD apresenta-se como uma abordagem relevante para reduzir a dependência de palavras de ativação explícitas e promover interações mais fluidas e intuitivas \cite{Wagner2024}.

Do ponto de vista acadêmico, o estudo de métodos alternativos de ativação baseados na detecção de intenção implícita atua na interseção de três áreas de pesquisa: (i) reconhecimento de fala, responsável pela transcrição do sinal acústico; (ii) processamento de linguagem natural, voltado à interpretação semântica; e (iii) interação humano-computador, que modela os aspectos de interação do sistema. Estudos recentes indicam que abordagens multimodais para DDSD têm alcançado melhorias consistentes em métricas de desempenho em comparação a modelos unimodais de texto ou áudio \cite{Wagner2024}. Entretanto, essas abordagens geralmente demandam alto poder computacional, grandes volumes de dados multimodais rotulados e processos complexos de sincronização entre modalidades, o que impõe barreiras significativas para o treinamento e a aplicação em dispositivos com recursos restritos \cite{Xu2024}.

Além disso, observa-se uma lacuna relevante na literatura: os trabalhos sobre DDSD foram desenvolvidos predominantemente para o idioma inglês, empregando abordagens multimodais ou puramente acústicas \cite{Norouzian2019, Krishna2023}. O português brasileiro, apesar de ser um dos idiomas mais falados do mundo, permanece sub-representado nessa área, com escassez de estudos voltados à detecção de intenção implícita. 

Embora haja avanços recentes na disponibilização de recursos linguísticos para o português brasileiro, como modelos pré-treinados, a exemplo do BERTimbau \cite{Souza2020}, e \textit{datasets} de fala, como o CORAA ASR \cite{CORAA2023}, o NURC-SP \cite{NURCSP2025} e o TAGARELA \cite{Tagarela2026}, esses recursos são majoritariamente voltados para tarefas de reconhecimento de fala. Observa-se, portanto, uma escassez de corpora especificamente anotados para DDSD em português brasileiro, o que limita o desenvolvimento e a avaliação de métodos voltados à detecção de intenção implícita nesse idioma.

Diante desse cenário, este trabalho busca contribuir para o preenchimento dessa lacuna ao investigar uma abordagem de DDSD baseada exclusivamente no nível textual, a partir de transcrições geradas por sistemas ASR e utilizando modelos de linguagem pré-treinados para o português brasileiro. A escolha da modalidade textual justifica-se por três fatores principais: (i) modelos baseados em arquitetura Transformer alcançam estado da arte em tarefas de classificação de intenção e inferência semântica \cite{Weld2022}; (ii) experimentos recentes demonstram que aplicar modelos de linguagem sobre hipóteses de sistemas ASR fornece um sinal discriminativo relevante para distinguir entre fala dirigida e não dirigida ao sistema \cite{Wagner2024, Dighe2024}; e (iii) a abordagem em nível de texto garante maior modularidade arquitetural, facilitando a comparação com \textit{baselines} clássicos e a análise do impacto de erros de transcrição, configurando-se como uma alternativa prática e escalável \cite{Dighe2024}.

Ainda que abordagens multimodais frequentemente alcancem melhor desempenho em métricas como a Taxa de Erro Igual, em inglês \textit{Equal Error Rate} (EER) \cite{Wagner2024}, a simplicidade, interpretabilidade e viabilidade prática da abordagem puramente textual configuram-se como uma alternativa relevante no contexto de recursos limitados para o português brasileiro. A abordagem textual não é apenas uma limitação de escopo, mas um decisão estratégica: permite criar um \textit{baseline} inicial sólido que, no futuro, pode ser combinado com informações acústicas e contextuais. O desenvolvimento deste sistema atua como um \textit{baseline} inicial para subsidiar pesquisas futuras em DDSD no contexto específico do português brasileiro.

Vale ressaltar também que sistemas que detectam intenção de fala em ambiente contínuo levantam questões relevantes de privacidade e ética: escuta ambiente, ativação indevida, armazenamento de áudio e proteção de dados do usuário. Este trabalho não aborda implantação em dispositivos reais nem armazenamento contínuo de áudio, concentrando-se exclusivamente na classificação textual de transcrições em ambiente controlado. Estudos futuros devem considerar políticas de privacidade, processamento local (\textit{on-device}) e mecanismos de consentimento do usuário \cite{Chen2021}.

% ---------------------------------------------------------------
\section{Objetivos}
\label{sec:objetivos}


\subsection{Objetivo principal}

O objetivo principal deste trabalho é avaliar a viabilidade e o desempenho de modelos de linguagem pré-treinados para o português brasileiro na tarefa de Detecção de Fala Direcionada ao Dispositivo (DDSD), operando exclusivamente sobre transcrições textuais geradas por sistemas ASR, e verificar se essa abordagem alcança desempenho superior a \textit{baselines} textuais tradicionais (TF-IDF + SVM/Regressão Logística) e desempenho competitivo em relação a abordagens unimodais reportadas na literatura, em termos de F1-score, acurácia e EER.

\subsection{Objetivos específicos}

Para que o objetivo principal seja alcançado, foram definidos os seguintes objetivos específicos:

\begin{enumerate}
    \item Construir um corpus de sentenças em português brasileiro anotadas quanto à
    direção da fala (dirigida ao dispositivo \textit{vs.} conversa incidental), a partir
    de adaptações de \textit{datasets} existentes (CORAA ASR \cite{CORAA2023},
    NURC-SP \cite{NURCSP2025}, TAGARELA \cite{Tagarela2026}) e de sentenças
    sintetizadas em ambiente controlado, com critérios explícitos de anotação, protocolo
    de concordância entre anotadores e divisão balanceada entre treino, validação e
    teste;

    \item Implementar e avaliar um classificador binário baseado no BERTimbau
    \cite{Souza2020} como modelo principal, com BERTimbau ajustado por
    \textit{fine-tuning} para a tarefa de DDSD, e comparar seu desempenho com um
    \textit{baseline} clássico baseado em TF-IDF combinado com SVM ou Regressão
    Logística \cite{Manning2008}, utilizando como entrada as transcrições geradas pelo
    Whisper \cite{Radford2023};

    \item Comparar o desempenho do classificador proposto com baselines estabelecidos na
    literatura de DDSD, utilizando métricas padronizadas — F1-score, acurácia, precisão,
    \textit{recall} e EER — e discutir os resultados em relação a abordagens unimodais
    reportadas na literatura para o inglês, deixando explícito que essa comparação tem
    caráter aproximado, dado que os \textit{datasets} e as modalidades diferem
    \cite{Wagner2024};

    \item Analisar o impacto de erros de transcrição do sistema ASR no desempenho do
    classificador de intenção, comparando os resultados obtidos com transcrições manuais
    e com transcrições automáticas, e medindo a correlação entre a Taxa de Erro de Palavras
    (\textit{Word Error Rate}, WER) do ASR e a degradação do desempenho do classificador
    \cite{Dighe2024}.
\end{enumerate}

% ---------------------------------------------------------------
\section{Hipótese de pesquisa}

A hipótese central deste trabalho é a seguinte:

\textit{Modelos de linguagem pré-treinados para o português brasileiro, aplicados sobre
transcrições geradas por sistemas ASR, conseguem alcançar desempenho superior a
\textit{baselines} textuais tradicionais (TF-IDF + SVM ou Regressão Logística) e
desempenho competitivo em relação a abordagens acústicas unimodais reportadas na
literatura, considerando métricas como F1-score, acurácia e EER avaliadas sobre um
corpus anotado em português brasileiro.}

Essa hipótese é fundamentada em três evidências da literatura. Primeiro, estudos recentes
demonstram que a análise textual de hipóteses ASR fornece um sinal discriminativo
relevante e complementar ao sinal acústico para a tarefa de DDSD \cite{Wagner2024,
Dighe2024}. Segundo, modelos baseados na arquitetura Transformer alcançam estado da arte
em tarefas de classificação de intenção textual, superando abordagens tradicionais de
aprendizado de máquina \cite{Weld2022}. Terceiro, modelos pré-treinados especificamente
para o português brasileiro, como o BERTimbau \cite{Souza2020} e o DeBERTinha
\cite{Abonizio2023}, têm demonstrado desempenho competitivo em tarefas de classificação
textual, o que fundamenta sua aplicabilidade à tarefa proposta.

A hipótese será testada por meio de experimentos controlados no TCC2, com avaliação
quantitativa sobre um corpus anotado em português brasileiro, utilizando protocolo de
validação cruzada e comparação explícita com os \textit{baselines} definidos nos
Objetivos Específicos.

\section{Limitações do trabalho}


Este trabalho apresenta limitações inerentes ao seu escopo e ao contexto em que será
desenvolvido.

A primeira limitação diz respeito à modalidade de entrada. A abordagem proposta opera
exclusivamente sobre o nível textual, descartando pistas acústicas como prosódia,
entonação e características paralinguísticas que são reconhecidamente úteis para a
tarefa de DDSD \cite{Norouzian2019, Krishna2023}. Essa é uma escolha estratégica de
escopo: o objetivo é criar um \textit{baseline} textual sólido antes de incorporar
modalidades adicionais. Ainda assim, implica que o sistema pode apresentar desempenho
inferior a abordagens multimodais em cenários com alta variabilidade acústica.

A segunda limitação refere-se ao impacto dos erros de transcrição. A qualidade das
saídas geradas pelo sistema ASR afeta diretamente o classificador textual: erros de
transcrição podem degradar o desempenho do modelo, especialmente em falas com sotaque,
ruído de fundo ou vocabulário incomum \cite{Bevilacqua2024}. O impacto desses erros será
analisado como parte dos experimentos, mas representa uma fonte de variabilidade não
completamente controlável. A comparação entre transcrições manuais e automáticas, prevista
no quarto objetivo específico, buscará quantificar essa degradação.

A terceira limitação envolve a construção do corpus. A escassez de corpora
especificamente anotados para DDSD em português brasileiro obriga a construção de um
corpus próprio em ambiente controlado. Isso implica limitações em termos de escala,
diversidade dialetal e espontaneidade. O corpus resultante não cobrirá toda a variabilidade
da fala espontânea em contextos reais \cite{CORAA2023}, o que pode limitar a generalização
dos resultados.

Por fim, este trabalho não aborda aspectos de privacidade, armazenamento de áudio ou
implantação contínua em dispositivos reais, concentrando-se exclusivamente na
classificação textual de transcrições em ambiente controlado. Estudos futuros devem
considerar políticas de privacidade, processamento local e mecanismos de consentimento
do usuário \cite{Chen2021}.

% ===============================================================
% CAPÍTULO 2 - REVISÃO BIBLIOGRÁFICA
% ===============================================================
\chapter{Revisão Bibliográfica}
\label{cap:revisao}

% ---------------------------------------------------------------
\section{Referencial teórico}
\label{sec:referencial}


% ---------------------------------------------------------------
\section{Metodologia da pesquisa bibliográfica}
\label{sec:metodologia_pesquisa}

A pesquisa bibliográfica foi conduzida com o objetivo de identificar trabalhos científicos
relevantes relacionados à detecção de intenção em linguagem natural, assistentes virtuais
baseados em voz e técnicas de aprendizado de máquina aplicadas ao reconhecimento de
comandos de voz.
Foram consultadas bases de dados acadêmicas amplamente utilizadas na área de
Computação, incluindo Google Scholar, Semantic Scholar e arXiv. Essas plataformas foram
selecionadas por concentrarem artigos científicos, pré-publicações e trabalhos revisados
por pares.
As buscas foram realizadas utilizando combinações de palavras-chave em português
e inglês, tais como: “voice assistant”, “intent detection”, “keyword spotting”, “implicit
intent”, “natural language processing” e “voice command classification”. Essas expressões
foram combinadas utilizando operadores booleanos, como AND e OR, com o objetivo de
ampliar ou restringir os resultados conforme necessário.
Como critérios de inclusão, foram considerados trabalhos publicados entre os anos de
2021 e 2026, disponíveis nos idiomas inglês ou português, e que apresentassem relação
direta com o problema de pesquisa abordado neste trabalho. Foram priorizados artigos
que apresentassem experimentos, métodos de classificação ou avaliações de sistemas de
interação por voz.
Inicialmente, foi realizada uma triagem baseada na leitura do título e do resumo dos
artigos encontrados. Trabalhos considerados relevantes foram posteriormente analisados
com maior profundidade para compor a fundamentação teórica e a seção de trabalhos
relacionados deste estudo.

% ---------------------------------------------------------------
\section{Fundamentos}
\label{sec:fundamentos}


% ---------------------------------------------------------------
\section{Trabalhos relacionados}
\label{sec:trabalhos_relacionados}



% ===============================================================
% CAPÍTULO 3 - DESENVOLVIMENTO
% ===============================================================
\chapter{Desenvolvimento}
\label{cap:desenvolvimento}

% ---------------------------------------------------------------
\section{Solução proposta}
\label{sec:solucao}


% ---------------------------------------------------------------
\section{Cronograma de trabalho}
\label{sec:cronograma}


% ---------------------------------------------------------------
% ELEMENTOS PÓS-TEXTUAIS
% ---------------------------------------------------------------
\postextual

% ===============================================================
% CAPÍTULO 4 - REFERÊNCIAS BIBLIOGRÁFICAS
% ===============================================================
% As referências bibliográficas são geradas automaticamente a partir
% do arquivo bibliografia.bib, seguindo o padrão ABNT (NBR 6023).
%
% COMO CITAR no texto:
%   \cite{chave}         -> (SOBRENOME, ano) — citação indireta
%   \citeonline{chave}   -> Sobrenome (ano) — citação direta no texto
%
% TIPOS DE REFERÊNCIA no arquivo .bib:
%   @book        -> Livros
%   @article     -> Artigos de periódicos
%   @inproceedings -> Artigos em anais de eventos
%   @mastersthesis -> Dissertações de mestrado
%   @phdthesis   -> Teses de doutorado
%   @misc        -> Sites, vídeos e outros materiais
%   @manual      -> Documentação técnica
%
% Veja o arquivo bibliografia.bib para exemplos de cada tipo.
% ---------------------------------------------------------------
\bibliography{bibliografia}

% ---------------------------------------------------------------
% APÊNDICES (se necessário)
% ---------------------------------------------------------------
% \begin{apendicesenv}
% \partapendices
%
% \chapter{Nome do Apêndice}
% \label{ap:apendice1}
%
% Conteúdo do apêndice (material elaborado pelo próprio autor).
%
% \end{apendicesenv}

% ---------------------------------------------------------------
% ANEXOS (se necessário)
% ---------------------------------------------------------------
% \begin{anexosenv}
% \partanexos
%
% \chapter{Nome do Anexo}
% \label{an:anexo1}
%
% Conteúdo do anexo (material de terceiros).
%
% \end{anexosenv}

\end{document}
